%% VocalCoach_Kim — Detailed Technical Report
%% Compile: pdflatex report.tex  (run twice for TOC/references)
\documentclass[11pt,letterpaper]{article}

\usepackage{amsmath,amssymb}
\usepackage{booktabs}
\usepackage{array}
\usepackage{xcolor}
\usepackage{geometry}
\usepackage{hyperref}
\usepackage{parskip}
\usepackage{enumitem}
\usepackage{fancyhdr}
\usepackage{titlesec}
\usepackage{listings}
\usepackage{caption}
\usepackage{subcaption}
\usepackage{microtype}

\geometry{margin=1in}

\hypersetup{
  colorlinks=true,
  linkcolor=blue!60!black,
  urlcolor=blue!60!black,
  citecolor=blue!60!black
}

\definecolor{codegray}{RGB}{245,245,245}
\definecolor{codeframe}{RGB}{200,200,200}

\lstset{
  backgroundcolor=\color{codegray},
  basicstyle=\ttfamily\small,
  frame=single,
  rulecolor=\color{codeframe},
  breaklines=true,
  columns=flexible,
  keepspaces=true,
  showstringspaces=false
}

\pagestyle{fancy}
\fancyhf{}
\rhead{VocalCoach\_Kim Technical Report}
\lhead{IMAPLE Lab --- Drexel University}
\rfoot{\thepage}

\titleformat{\section}{\large\bfseries}{}{0em}{\thesection\quad}
\titleformat{\subsection}{\normalsize\bfseries}{}{0em}{\thesubsection\quad}

%% ── Notation shortcuts ────────────────────────────────────────────────────────
\newcommand{\fh}{\hat{f}}
\newcommand{\vh}{\hat{v}}
\newcommand{\sh}{\hat{s}}
\newcommand{\eh}{\hat{e}}
\newcommand{\ah}{\hat{a}}
\newcommand{\bh}{\hat{b}}
\newcommand{\dh}{\hat{d}}

% ─────────────────────────────────────────────────────────────────────────────
\begin{document}

\begin{titlepage}
  \centering
  \vspace*{2cm}
  {\LARGE\bfseries Voice Coaching System\\[6pt]
   for Singing Performance Evaluation\par}
  \vspace{0.5cm}
  {\large\itshape Detailed Technical Report\par}
  \vspace{1.5cm}
  {\large IMAPLE Lab, Drexel University\par}
  \vspace{0.5cm}
  {\large 2026\par}
  \vspace{2cm}
  \begin{abstract}
    \noindent
    This report provides a comprehensive technical description of the
    \texttt{VocalCoach\_Kim} software system — a reference-guided pipeline for
    objective singing performance evaluation.
    The system accepts a raw audio recording and a musical reference (MusicXML
    score and Praat TextGrid annotation) and produces a structured coaching
    report covering pitch accuracy, timing, note duration, and lyric
    intelligibility.
    Three deep neural networks act as \emph{feature extractors}: a
    torchcrepe/CREPE CNN for fundamental frequency estimation, a fine-tuned
    Wav2Vec2 model with a three-class linear head for note onset/offset
    detection, and a pre-trained Wav2Vec2 CTC model for phoneme boundary
    detection.
    All coaching scores are derived from explicit mathematical equations applied
    to extracted parameters; no model directly predicts a score, making the
    system fully interpretable and deterministic.
    This report describes every pipeline stage in detail, including data
    representations, algorithmic choices, configuration parameters, and the
    mathematical formulae underlying all metrics and scores.
  \end{abstract}
\end{titlepage}

\tableofcontents
\newpage

% =============================================================================
\section{System Overview}
% =============================================================================

The \texttt{VocalCoach\_Kim} system is organized as a seven-phase sequential
pipeline orchestrated by \texttt{UnifiedInferencePipeline}
(\texttt{inference/pipeline.py}).
Each phase is described in detail in subsequent sections; here we give a
high-level map.

\begin{table}[h]
  \centering
  \caption{Pipeline phases and responsible modules.}
  \renewcommand{\arraystretch}{1.25}
  \begin{tabular}{@{}clp{9cm}@{}}
    \toprule
    Phase & Module & Description \\
    \midrule
    1 & \texttt{utils/audio.py}
      & Load and normalize audio to 16\,kHz mono float32. \\
    2 & \texttt{models/pitch/pipeline.py}
      & Extract frame-level F0 and voiced/unvoiced mask via VAD + torchcrepe. \\
    3 & \texttt{models/phoneme/phoneme\_model.py}
      & Predict phoneme segment boundaries and labels via Wav2Vec2 + CTC. \\
    4 & \texttt{models/onset\_offset/model\_v5.py}
      & Predict note onset and offset times via fine-tuned Wav2Vec2. \\
    5 & \texttt{fusion/}
      & Resample all outputs to a canonical 100\,fps grid; construct
        \texttt{NoteEvent}, \texttt{LyricEvent}, and \texttt{WordEvent}
        structures. \\
    6 & \texttt{reference/} + \texttt{alignment/}
      & Parse MusicXML and TextGrid reference files; greedily align predicted
        events against reference events. \\
    7 & \texttt{metrics/} + \texttt{scoring/}
      & Compute objective metrics; normalize to [0,\,100]; generate final
        score report and rule-based interpretation. \\
    \bottomrule
  \end{tabular}
\end{table}

\textbf{Design principle.}
Neural models estimate measurable physical parameters
$(\fh_t,\,\vh_t,\,\sh_n,\,\eh_n,\,\ah_p,\,\bh_p)$.
All coaching intelligence is encoded in deterministic equations that map those
parameters to scores.
This separation makes the output of every computation auditable, reproducible,
and independent of training-time distributional assumptions.

\textbf{Canonical frame grid.}
All frame-level data is ultimately aligned to a grid with hop length
$H = 160$ samples at $f_s = 16\,000$\,Hz, giving a frame duration of
$H/f_s = 10$\,ms and a frame rate of 100\,fps.
The total number of canonical frames for a recording of $T$ samples is
$N_\text{frames} = \lceil T / H \rceil$.

\textbf{Dataset.}
The system was developed and evaluated using the GTSinger dataset, English
alto singer subset (\texttt{EN-Alto-1}).

% =============================================================================
\section{Notation}
% =============================================================================

Table~\ref{tab:notation} defines all symbols used throughout this report.
A hat ($\hat{\,\cdot\,}$) always denotes a quantity \emph{predicted} from
audio; unhatted symbols denote \emph{reference} quantities from score files.

\begin{table}[h]
  \centering
  \caption{Notation used throughout the report.}
  \label{tab:notation}
  \renewcommand{\arraystretch}{1.3}
  \begin{tabular}{@{}lp{11cm}@{}}
    \toprule
    Symbol & Meaning \\
    \midrule
    $\fh_t$                & Predicted fundamental frequency (Hz) at frame $t$ \\
    $\vh_t \in \{0,1\}$   & Predicted voiced/unvoiced flag at frame $t$ \\
    $r_t$                  & Reference F0 (Hz) at frame $t$ (derived from MIDI pitch) \\
    $\sh_n,\,\eh_n$        & Predicted onset and offset time (s) of note $n$ \\
    $\dh_n = \eh_n-\sh_n$ & Predicted duration of note $n$ \\
    $s_n,\,e_n,\,d_n$     & Reference onset, offset, and duration of note $n$ \\
    $m_n$                  & MIDI pitch of reference note $n$ \\
    $r_n = 440\cdot 2^{(m_n-69)/12}$ & Reference frequency (Hz) of note $n$ \\
    $\delta_n$             & Pitch deviation in cents for matched note $n$,
                             $1200\log_2(\text{median}\,\fh / r_n)$ \\
    $\epsilon_n$           & Onset deviation (ms), $(\sh_n - s_n)\times 1000$ \\
    $\varepsilon_n$        & Offset deviation (ms), $(\eh_n - e_n)\times 1000$ \\
    $\ah_p,\,\bh_p$        & Predicted phoneme $p$ start and end time (s) \\
    $a_p,\,b_p$            & Reference phoneme $p$ start and end time (s) \\
    $V$                    & Set of frames where both $\fh_t > 0$ and $r_t > 0$ \\
    $N$                    & Number of matched note pairs \\
    $P$                    & Number of matched phoneme pairs \\
    $\tau$                 & Tolerance threshold (cents or ms, context-dependent) \\
    \bottomrule
  \end{tabular}
\end{table}

\textbf{Sign conventions.}
A positive onset deviation $\epsilon_n > 0$ means the singer entered the note
\emph{later} than the reference.
A positive cent error $\delta_n > 0$ means the singer was \emph{sharp}
relative to the reference pitch.

% =============================================================================
\section{Audio Preprocessing (Phase 1)}
% =============================================================================

The entry point for audio loading is \texttt{utils/audio.py}.
All incoming audio is converted to a canonical representation before any model
sees it:

\begin{itemize}[noitemsep]
  \item Sample rate: $f_s = 16\,000$\,Hz (resampled if necessary).
  \item Channels: mono (multi-channel audio is averaged across channels).
  \item Dtype: 32-bit float, normalized to the range $[-1, +1]$.
  \item Frame grid: hop length $H = 160$ samples, giving 10\,ms frames.
\end{itemize}

The normalization step applies peak normalization so that the loudest sample
in the file has magnitude 1.0.
All downstream models operate on this canonical waveform array.

% =============================================================================
\section{Pitch Extraction (Phase 2)}
% =============================================================================

Pitch extraction is handled by \texttt{models/pitch/pipeline.py} through the
class \texttt{PitchVADPipeline}.
The pipeline has five stages.

\subsection{Voice Activity Detection}

\texttt{models/pitch/vad.py} implements \texttt{WebRTCVAD}, a wrapper around
the WebRTC voice activity detector library.
Key parameters:
\begin{itemize}[noitemsep]
  \item \textbf{Aggressiveness}: mode 2 (0 = least aggressive, 3 = most).
  \item \textbf{Frame duration}: 20\,ms.
  \item \textbf{Smoothing}: median filter of window size 5 applied to the raw
        binary mask to remove isolated flip decisions.
  \item \textbf{Fallback}: if the WebRTC library is not installed, an
        energy-based detector activates with a threshold of $-40$\,dBFS.
\end{itemize}
Output: a boolean array \texttt{voiced\_mask} of shape $(N_\text{vad},)$,
where $N_\text{vad}$ is the number of 20\,ms VAD frames in the recording,
together with corresponding frame-center timestamps.

\subsection{Fundamental Frequency Estimation}

\texttt{models/pitch/pitch\_wrapper.py} implements \texttt{PitchModelWrapper},
which supports three backends:

\textbf{torchcrepe (primary).}
A CREPE-style convolutional neural network~\cite{kim2018crepe} trained to
predict a probability distribution over 360 frequency bins spanning
32.7--1975.5\,Hz.
Configuration:
\begin{itemize}[noitemsep]
  \item Hop length: 160 samples (10\,ms at 16\,kHz).
  \item Frequency range: 50--1000\,Hz.
  \item Model capacity: \texttt{full} (6-layer CNN).
  \item Decoder: Viterbi path decoding for temporal smoothness.
  \item Periodicity threshold: 0.21 (frames with confidence below this are
        treated as unvoiced).
  \item Silence threshold: $-60$\,dBFS.
  \item Batch size: 1024.
\end{itemize}
Output: \texttt{(times, f0, confidence)} all as float32 arrays of length
$N_\text{frames}$.

\textbf{pYIN (fallback).}
If torchcrepe is unavailable, the system falls back to librosa's pYIN
implementation with \texttt{fmin}=65\,Hz, \texttt{fmax}=2093\,Hz, frame
length 2048, and hop 512.

\textbf{Custom backend.}
Users may supply any callable returning a \texttt{(times, f0, confidence)}
triple.

After prediction, all NaN values are replaced with 0.0, and the output is
cast to float32.

\subsection{VAD Alignment}

The VAD mask is defined at a 50\,Hz frame rate (20\,ms frames), while the
pitch output is at 100\,fps (10\,ms frames).
The mask is resampled to the pitch frame rate using nearest-neighbour
interpolation:
\[
  \vh_t^\text{vad} = \texttt{voiced\_mask}\!\left[\left\lfloor
      \frac{t \cdot f_\text{vad}}{f_\text{pitch}}
    \right\rfloor\right]
\]
where $f_\text{vad}=50$ and $f_\text{pitch}=100$.

\subsection{VAD--Pitch Fusion}

\texttt{models/pitch/fusion.py} implements \texttt{fuse\_vad\_and\_pitch()},
which performs five sequential operations to produce a clean F0 trajectory.

\begin{enumerate}[noitemsep]
  \item \textbf{Intersection.}
        $\vh_t^\text{final} \leftarrow \vh_t^\text{pitch} \wedge
        \vh_t^\text{vad}$.
        A frame is voiced only if both the pitch model and the VAD agree.
  \item \textbf{Masking.}
        Unvoiced frames are zeroed: $\fh_t \leftarrow 0$ whenever
        $\neg\vh_t^\text{final}$.
  \item \textbf{Gap interpolation.}
        Contiguous unvoiced gaps of up to 10 frames ($\approx 100$\,ms) are
        linearly interpolated between the F0 values on either boundary, and
        the corresponding frames are marked voiced.
        This prevents short closure events (e.g., consonants in a legato
        phrase) from fragmenting the pitch contour.
  \item \textbf{Median filtering.}
        A median filter of kernel size 5 is applied to voiced frames only.
        This suppresses isolated pitch estimation errors without smearing
        formant transitions.
  \item \textbf{Gaussian smoothing.}
        A Gaussian kernel with $\sigma = 1.0$ frame is convolved with each
        contiguous voiced segment.
\end{enumerate}

\textbf{Output.}
The pipeline returns \texttt{PipelineOutput(times, f0, voiced\_mask)} at
100\,fps.
The three arrays have shape $(N_\text{frames},)$ and dtype float32 / bool.

% =============================================================================
\section{Note Onset/Offset Detection (Phase 4)}
% =============================================================================

\subsection{Model Architecture}

The onset/offset model (\texttt{models/onset\_offset/model\_v5.py}) is built on
\texttt{facebook/wav2vec2-base}.
Wav2Vec2 consists of two subnetworks:

\textbf{Feature extractor (frozen).}
A stack of seven 1-D convolutional layers with a total stride of 320 samples.
At 16\,kHz, this maps the raw waveform to a feature sequence at
$16\,000/320 = 50$\,Hz (20\,ms per frame).
These layers are frozen during fine-tuning (parameters are not updated).

\textbf{Transformer encoder (fine-tuned).}
Twelve transformer layers with hidden dimension $d = 768$.
These layers are fine-tuned on the singing dataset.

\textbf{Dropout and head.}
A dropout layer with $p = 0.1$ followed by a single linear projection:
\[
  (\ell_t^{(0)},\,\ell_t^{(1)},\,\ell_t^{(2)}) = \mathbf{W}\,\mathbf{h}_t + \mathbf{b}
  \quad \in \mathbb{R}^3
\]
where $\mathbf{h}_t \in \mathbb{R}^{768}$ is the transformer output at frame
$t$.
The three output dimensions correspond to onset, offset, and active (voiced)
probabilities respectively.

\subsection{Inference and Peak Picking}

Logits are passed through a sigmoid:
\[
  p_t^{(k)} = \sigma\!\left(\ell_t^{(k)}\right) \in [0,1],
  \quad k \in \{0_\text{onset},\, 1_\text{offset},\, 2_\text{active}\}
\]

Peak picking uses \texttt{scipy.signal.find\_peaks} on the onset and offset
probability curves independently:
\[
  \sh_n = \mathrm{PeakPick}\!\left(p_t^{(0)},\;
    \theta_\text{on}{=}0.3,\; d_\text{min}{=}3\,\text{frames}\right)
\]
\[
  \eh_n = \mathrm{PeakPick}\!\left(p_t^{(1)},\;
    \theta_\text{off}{=}0.3,\; d_\text{min}{=}3\,\text{frames}\right)
\]
The minimum peak distance of 3 frames corresponds to approximately 60\,ms,
which prevents a single note boundary from producing multiple spurious
detections.

Onset--offset pairing is greedy and left-to-right: each onset is matched to
the next available offset that occurs strictly after it.
If no offset is found, a synthetic offset is inserted 10\,ms before the next
onset (or 300\,ms after the last onset for the final note).

The model runs at 50\,Hz; outputs are resampled to the canonical 100\,fps
grid in the fusion step (Section~\ref{sec:fusion}).

\subsection{Training Configuration}

The model is fine-tuned with the following settings
(\texttt{configs/onset\_offset.yaml}):
batch size 16, 100 epochs, Adam optimizer with learning rate $10^{-3}$,
patience 15 epochs (early stopping), and a Gaussian label sigma of 20\,ms
(the binary onset/offset ground-truth annotations are soft-labeled with a
Gaussian to account for annotation uncertainty).

% =============================================================================
\section{Phoneme Boundary Detection (Phase 3)}
% =============================================================================

\subsection{Model}

The phoneme model (\texttt{models/phoneme/phoneme\_model.py}) uses
\texttt{facebook/wav2vec2-lv-60-espeak-cv-ft}, a Wav2Vec2 model pre-trained on
60 languages and fine-tuned for CTC phoneme recognition using the eSpeak IPA
vocabulary.
Unlike the onset/offset model, this model is used \emph{purely for inference}
--- no layers are frozen or fine-tuned; it is loaded in evaluation mode
(\texttt{model.eval()}) and weights remain fixed.

\subsection{CTC Alignment Pipeline}

\begin{enumerate}[noitemsep]
  \item \textbf{Feature extraction.}
        The raw 16\,kHz waveform is passed through the Wav2Vec2 feature
        extractor (frozen CNN), then through the transformer encoder to
        produce a per-frame logit tensor
        $L \in \mathbb{R}^{T \times |V|}$ where $|V|$ is the CTC vocabulary
        size (includes blank token, ID = 0).
  \item \textbf{CTC decoding.}
        Argmax decoding: $\hat{y}_t = \arg\max_v L_{t,v}$.
        Consecutive repeated tokens are collapsed.
        Blank tokens (ID = 0) are removed.
        The result is a sequence of (token\_id, start\_frame, end\_frame)
        triples.
  \item \textbf{Time conversion.}
        Frame indices are converted to seconds using the frame duration of
        the Wav2Vec2 CNN ($\approx$20--40\,ms depending on the specific layer
        stride).
  \item \textbf{Confidence scoring.}
        For each segment, the mean of the softmax probability
        $\mathrm{softmax}(L_{t,\cdot})[\text{token\_id}]$ over frames in
        the segment window is used as a confidence score.
\end{enumerate}

\subsection{Post-Processing}

Two post-processing steps recover phonemes that CTC alignment tends to
suppress.

\textbf{Over-long segment splitting.}
Any segment longer than 300\,ms is split.
The split point is identified as the interior frame (excluding the first and
last 10\,ms) where an alternative phoneme achieves the highest probability
above a threshold of 0.25.
This handles long sustained vowels that CTC assigns as a single segment
spanning the entire duration.

\textbf{Blank-region scanning.}
After the initial CTC output, the inter-segment gaps (blank regions) are
scanned.
For each gap longer than one frame, the frame-level softmax probabilities are
examined.
If any non-blank token achieves a peak probability above 0.05, that phoneme
is inserted as a new segment at the peak frame.
This recovers short consonants (e.g., plosives, fricatives) that CTC's
implicit blank-dominance tends to skip during singing, where consonants are
often de-emphasized.

\textbf{Output.}
A list of \texttt{PhonemeSegment} objects:
\begin{lstlisting}
PhonemeSegment(phoneme: str, start_time: float,
               end_time: float, confidence: float)
\end{lstlisting}
Phoneme labels are normalized to IPA using
\texttt{reference/phoneme\_normalization.py}.

% =============================================================================
\section{Fusion and Canonical Representation (Phase 5)}
\label{sec:fusion}
% =============================================================================

\subsection{Frame-Rate Alignment}

\texttt{fusion/alignment.py} implements \texttt{merge\_model\_outputs()},
which resamples all model outputs to the canonical 100\,fps grid.

\begin{itemize}[noitemsep]
  \item \textbf{Float arrays} (F0, onset/offset probability curves):
        linearly interpolated to the canonical timestamps;
        out-of-range values filled with 0.0.
  \item \textbf{Boolean masks} (voiced flags):
        nearest-neighbour resampling; a value $> 0.5$ maps to \texttt{True}.
  \item \textbf{Phoneme segment list} to per-frame labels:
        each canonical frame is assigned the label of the phoneme segment
        whose interval has the largest overlap with that frame's window.
\end{itemize}

Output: \texttt{FrameAlignedFeatures}
\begin{lstlisting}
FrameAlignedFeatures(
    timestamps:     ndarray (N_frames,)   # seconds
    f0:             ndarray (N_frames,)   # Hz, 0 if unvoiced
    voiced:         ndarray (N_frames,)   # bool
    onset_probs:    ndarray (N_frames,)   # [0,1]
    offset_probs:   ndarray (N_frames,)   # [0,1]
    phoneme_labels: List[str] length N_frames
)
\end{lstlisting}

\subsection{Note Event Construction}

\texttt{fusion/note\_events.py} constructs \texttt{NoteEvent} objects from the
aligned onset/offset probability curves.

\textbf{Peak picking} (on canonical 100\,fps grid):
\texttt{scipy.signal.find\_peaks} with threshold $\theta = 0.5$, minimum peak
distance corresponding to 50\,ms (\texttt{min\_distance\_s} from config).

\textbf{Greedy onset--offset pairing}: each onset is paired with the first
subsequent offset, subject to maximum note duration 10\,s.

\textbf{Pitch statistics} are computed per note from the fused F0 array.
For each note window $[\sh_n, \eh_n]$:
\begin{itemize}[noitemsep]
  \item Mean F0 of voiced frames in the window (used as the representative
        pitch).
  \item MIDI conversion: $\hat{m}_n = 69 + 12\log_2(\bar{f}_n/440)$.
  \item Pitch stability: standard deviation of voiced F0 values within the
        window.
  \item Voiced fraction: fraction of frames in the window with
        $\vh_t = \texttt{True}$.
\end{itemize}

\textbf{Note confidence}: $c_n = \sqrt{p_\text{onset} \cdot p_\text{offset}}$
where $p_\text{onset}$ and $p_\text{offset}$ are the peak probabilities at
the detected onset and offset frames respectively.

\subsection{Lyric and Word Event Construction}

\texttt{fusion/lyric\_events.py} converts the phoneme segment list into a
two-level hierarchy.

\textbf{LyricEvent} (one per phoneme): wraps \texttt{PhonemeSegment} and adds
a \texttt{word\_idx} and \texttt{note\_idx} back-reference.

\textbf{WordEvent}: consecutive phonemes whose start/end gap is below
\texttt{word\_gap\_s = 100\,ms} are grouped into a word.
The word text is the hyphen-joined sequence of its constituent phoneme strings.
Word confidence is the mean of constituent phoneme confidences.

\subsection{Cross-Event Annotation}

\texttt{fusion/event\_alignment.py} annotates notes with the phonemes that
overlap them (minimum overlap 10\,ms) and vice versa.
It also:
\begin{itemize}[noitemsep]
  \item Builds \emph{voiced regions}: contiguous spans of
        $\vh_t = \texttt{True}$.
  \item Builds \emph{phrase events}: groups of notes separated by inter-onset
        gaps below 500\,ms.
\end{itemize}

All of these structures are bundled into
\texttt{FusedPerformanceRepresentation}, the complete predicted-side
representation passed to phases 6 and 7.

% =============================================================================
\section{Reference Parsing (Phase 6a)}
% =============================================================================

\subsection{MusicXML Parsing}

\texttt{reference/musicxml\_parser.py} parses MusicXML files using the
\texttt{music21} library and extracts:

\begin{itemize}[noitemsep]
  \item \textbf{Tempo}: the first \texttt{MetronomeMark} object;
        default 120\,BPM if absent.
  \item \textbf{Time and key signatures}: first occurrences in the score.
  \item \textbf{Notes and rests}: obtained by flattening the score across all
        parts and iterating over all elements.
  \item \textbf{Tied notes}: consecutive tied notes are merged into a single
        event with combined duration.
\end{itemize}

\textbf{Time conversion.}
Musical beat positions are converted to absolute seconds:
\[
  t_\text{onset}(n) = \text{offset\_beats}(n) \times \frac{60}{\text{tempo\_bpm}}
\]

\textbf{Pitch conversion.}
MIDI pitch to frequency:
\[
  r_n = 440 \cdot 2^{(m_n - 69)/12}
\]

Output: \texttt{ReferencePerformanceRepresentation} containing a list of
\texttt{ReferenceNote} objects, each with fields: \texttt{onset\_time},
\texttt{offset\_time}, \texttt{duration}, \texttt{pitch\_midi},
\texttt{pitch\_hz}, \texttt{pitch\_name}, \texttt{lyric}, \texttt{measure},
\texttt{beat}, \texttt{is\_rest}, \texttt{is\_tied}, \texttt{note\_idx}.

\subsection{TextGrid Parsing}

\texttt{reference/textgrid\_parser.py} parses Praat TextGrid files using the
\texttt{praatio} library (with a plain-text fallback parser for edge cases).
It extracts two tiers:
\begin{itemize}[noitemsep]
  \item \textbf{Phoneme tier}: searched under names
        \texttt{phonemes}, \texttt{phones}, \texttt{phone}, \texttt{phoneme}.
  \item \textbf{Word tier}: searched under names
        \texttt{words}, \texttt{word}, \texttt{orthography}.
\end{itemize}

Silence intervals (labels in
$\{\texttt{""}, \texttt{"SIL"}, \texttt{"<SIL>"}, \texttt{"sp"},
   \texttt{"spn"}, \texttt{"sil"}, \texttt{"<eps>"}\}$)
are skipped by default.
Phoneme labels are normalized to IPA.
Each \texttt{ReferenceWord} stores indices of the phonemes that fall within
its time span.

\subsection{Reference Builder}

\texttt{reference/reference\_builder.py} combines the outputs of the two
parsers into the final \texttt{ReferencePerformanceRepresentation}.
It also builds phrase groupings from reference notes: a new phrase begins
whenever the gap between consecutive reference note onsets exceeds 500\,ms.

% =============================================================================
\section{Alignment (Phase 6b)}
% =============================================================================

\texttt{alignment/reference\_alignment.py} performs greedy one-to-one
matching between predicted and reference events for notes, phonemes, and words
separately.

\subsection{Note Alignment}

\begin{enumerate}[noitemsep]
  \item Reference rests are excluded (only pitched notes are matched).
  \item A temporal overlap matrix is computed:
        \[
          O_{ij} = \max\!\left(0,\;
            \min(\eh_i, e_j) - \max(\sh_i, s_j)
          \right)
        \]
        for all predicted note $i$ and reference note $j$ pairs.
  \item Pairs are sorted by descending overlap and matched greedily (each
        predicted and reference note participates in at most one match).
  \item A match is accepted only if
        $O_{ij} \ge 10$\,ms and
        $|\sh_i - s_j| \le 500$\,ms.
  \item For each accepted match, the following deviations are computed:
        \begin{align*}
          \text{onset deviation} &= (\sh_i - s_j) \times 1000 \;\text{ms}\\
          \text{offset deviation} &= (\eh_i - e_j) \times 1000 \;\text{ms}\\
          \delta_n &= 1200 \log_2\!\left(
            \frac{\text{median}\,\fh \text{ over }[\sh_i,\eh_i]}{r_j}
          \right) \;\text{cents}
        \end{align*}
\end{enumerate}

\subsection{Phoneme Alignment}

The same greedy overlap procedure is applied to phoneme pairs, with a
relaxed minimum overlap of 5\,ms.
Additionally, each match records a \texttt{label\_match} boolean
($\hat{l}_p = l_p$, predicted phoneme label equals reference label).

\subsection{Word Alignment}

Greedy overlap matching with minimum overlap 10\,ms and maximum onset
deviation 500\,ms.

\subsection{Output}

All matches are stored in \texttt{AlignmentResult}:
\begin{lstlisting}
AlignmentResult(
    note_matches:     List[NoteAlignmentMatch]
    phoneme_matches:  List[PhonemeAlignmentMatch]
    word_matches:     List[WordAlignmentMatch]
    unmatched_pred_notes:     List[int]
    unmatched_ref_notes:      List[int]
    unmatched_pred_phonemes:  List[int]
    ...
    alignment_metadata: Dict
)
\end{lstlisting}

% =============================================================================
\section{Metrics (Phase 7a)}
% =============================================================================

\subsection{Pitch Metrics}

All pitch metrics operate on the $N$ matched note pairs with valid pitch
deviation $\delta_n$ (in cents).

\textbf{Pitch Accuracy} ($\tau = 50$\,¢):
\[
  \mathrm{PitchAcc}_\tau = \frac{1}{N}\sum_{n=1}^{N}\mathbf{1}(|\delta_n| \le \tau)
\]

\textbf{Mean Absolute Cent Error (MACE)}:
\[
  \mathrm{MACE} = \frac{1}{N}\sum_{n=1}^{N}|\delta_n|
\]

\textbf{Root-Mean-Square Cent Error (RMSE)}:
\[
  \mathrm{PitchRMSE} = \sqrt{\frac{1}{N}\sum_{n=1}^{N}\delta_n^2}
\]

\textbf{Note-Level Pitch Accuracy} is numerically identical to
$\mathrm{PitchAcc}_\tau$ but populates a per-note breakdown field
(\texttt{PitchMetrics.per\_note}) used for note-by-note feedback.

All deviations $\delta_n$ come from the \texttt{pitch\_deviation\_cents} field
of each \texttt{NoteAlignmentMatch}, which is computed from the median F0 over
the predicted note window.

\subsection{Timing Metrics}

Let $\epsilon_n = (\sh_n - s_n) \times 1000$\,ms be the onset deviation and
$\varepsilon_n = (\eh_n - e_n) \times 1000$\,ms be the offset deviation.

\textbf{Onset Error Statistics}:
\[
  \bar{\epsilon} = \frac{1}{N}\sum_n\epsilon_n \quad\text{(mean bias)}
  \qquad
  \mathrm{MAE}_\text{on} = \frac{1}{N}\sum_n|\epsilon_n|
\]
\[
  \sigma_\text{on} = \sqrt{\frac{1}{N}\sum_n(\epsilon_n - \bar{\epsilon})^2}
  \quad\text{(consistency)}
\]

\textbf{Timing Accuracy} ($\tau = 50$\,ms):
\[
  \mathrm{TimingAcc}_\tau = \frac{1}{N}\sum_n\mathbf{1}(|\epsilon_n| \le \tau)
\]

\textbf{Offset Error}:
\[
  \mathrm{MAE}_\text{off} = \frac{1}{N}\sum_n|\varepsilon_n|
\]

\textbf{Inter-Onset Interval MAE (IOI-MAE)}.
The IOI sequence is $\Delta s^{(i)} = s_{i+1} - s_i$ for reference notes and
$\Delta\sh^{(i)} = \sh_{i+1} - \sh_i$ for predicted notes.
Let $M = \min(N_\text{pred}, N_\text{ref}) - 1$.
\[
  \mathrm{IOI\text{-}MAE} = \frac{1}{M}\sum_{i=1}^{M}
    |\Delta\sh^{(i)} - \Delta s^{(i)}| \times 1000\;\text{ms}
\]
IOI-MAE is agnostic to absolute timing offset: a singer who is consistently
200\,ms late but perfectly in rhythm will have IOI-MAE $\approx 0$ but high
$\mathrm{MAE}_\text{on}$.

\subsection{Duration Metrics}

Let $\dh_n = \eh_n - \sh_n$ (predicted) and $d_n = e_n - s_n$ (reference).

\textbf{Duration Error}:
\[
  \mathrm{DurErr} = \frac{1}{N}\sum_n|\dh_n - d_n| \;\text{s}
\]

\textbf{Relative Duration Error} (tempo-independent):
\[
  \mathrm{RelDurErr} = \frac{1}{N}\sum_n\frac{|\dh_n - d_n|}{d_n}
\]

\textbf{Duration Ratio}:
\[
  \mathrm{DurRatio} = \frac{1}{N}\sum_n\frac{\dh_n}{d_n}
\]
A ratio of 1.0 is perfect; $>1$ means the singer stretches notes; $<1$ means
notes are cut short.
For scoring, the deviation from unity $|\mathrm{DurRatio} - 1|$ is used.

\subsection{Lyric Metrics}

Let $P$ be the number of matched phoneme pairs.
Let $\epsilon_p^{(a)} = (\ah_p - a_p) \times 1000$\,ms be the onset
deviation for phoneme $p$.

\textbf{Phoneme Boundary Error}:
\[
  \mathrm{PhonBndErr} = \frac{1}{P}\sum_{p=1}^{P}|\epsilon_p^{(a)}|
\]
Only the start boundary is used (not the average of start and end).

\textbf{Word Alignment Accuracy}:
\[
  \mathrm{WordAcc} = \frac{|\text{matched words}|}{|\text{reference words}|}
\]

\textbf{Phoneme Overlap Accuracy}:
\[
  \mathrm{OverlapAcc} = \frac{1}{P}\sum_p
    \mathbf{1}\!\left(\frac{\min(\bh_p,b_p)-\max(\ah_p,a_p)}
                           {\bh_p - \ah_p} \ge 0.5\right)
\]

\textbf{Label Match Rate}:
\[
  \mathrm{LabelMatch} = \frac{1}{P}\sum_p\mathbf{1}(\hat{l}_p = l_p)
\]

% =============================================================================
\section{Scoring (Phase 7b)}
% =============================================================================

\subsection{Metric Normalization}

All raw metric values are mapped to the range $[0, 100]$ using one of two
normalization strategies.

\textbf{Piecewise-linear normalization} (for error metrics where lower is
better).
A set of breakpoints $\{(x_i, s_i)\}$ is defined such that $x_0 < x_1 <
\ldots$ and $s_0 > s_1 > \ldots$
(higher raw error maps to lower score).
For a raw value $x$ with $x_i \le x \le x_{i+1}$:
\[
  S_\text{pw}(x) = s_i + \frac{s_{i+1}-s_i}{x_{i+1}-x_i}(x - x_i)
\]
Values below $x_0$ are clamped to $s_0$; values above $x_{k}$ are clamped
to $s_k$.

\begin{table}[h]
  \centering
  \caption{Piecewise-linear breakpoints for error-based metrics.}
  \renewcommand{\arraystretch}{1.2}
  \begin{tabular}{@{}lcc@{}}
    \toprule
    Metric & $x$ breakpoints & $s$ breakpoints \\
    \midrule
    MACE (¢)            & 0, 25, 50, 100, 200   & 100, 88, 75, 50, 0 \\
    Pitch RMSE (¢)      & 0, 25, 50, 100, 200   & 100, 88, 72, 45, 0 \\
    Onset MAE (ms)      & 0, 25, 50, 100, 200   & 100, 88, 75, 50, 0 \\
    IOI MAE (ms)        & 0, 30, 60, 120, 240   & 100, 88, 75, 50, 0 \\
    Onset std (ms)      & 0, 30, 60, 120, 240   & 100, 88, 72, 45, 0 \\
    RelDurErr           & 0, 0.10, 0.20, 0.50, 1.0 & 100, 90, 75, 50, 0 \\
    Duration ratio dev  & 0, 0.10, 0.25, 0.50, 1.0 & 100, 90, 75, 50, 0 \\
    PhonBndErr (ms)     & 0, 15, 30, 60, 120    & 100, 88, 75, 50, 0 \\
    \bottomrule
  \end{tabular}
\end{table}

\textbf{Direct scaling} (for accuracy fractions where higher is better):
\[
  S_\text{acc}(a) = a \times 100, \quad a \in [0,1]
\]
Used for: PitchAcc, TimingAcc, WordAcc, OverlapAcc, LabelMatch.

\subsection{Confidence-Weighted Category Scores}

Each of the four category scores uses the same aggregation formula.
Let $\{(w_i, c_i, s_i)\}$ be the set of (weight, confidence, component score)
triples for the category.
Confidence is:
\[
  c_i = \min\!\left(1,\;\frac{N_i}{3}\right)
\]
where $N_i$ is the number of matched events contributing to component $i$.
This means fewer than 3 matched events yields partial confidence, ensuring
that data-poor evaluations contribute proportionally less.

\[
  S_\text{cat} = \frac{\displaystyle\sum_i w_i\,c_i\,s_i}
                      {\displaystyle\sum_i w_i\,c_i}
\]

\textbf{Pitch category} (weight 0.40 in final score):
\begin{center}
  \begin{tabular}{@{}lr@{}}
    \toprule
    Component & Weight $w_i$ \\
    \midrule
    $S_\text{acc}(\mathrm{PitchAcc})$ & 0.50 \\
    $S_\text{pw}(\mathrm{MACE})$      & 0.30 \\
    $S_\text{pw}(\mathrm{PitchRMSE})$ & 0.20 \\
    \bottomrule
  \end{tabular}
\end{center}

\textbf{Timing category} (weight 0.30 in final score):
\begin{center}
  \begin{tabular}{@{}lr@{}}
    \toprule
    Component & Weight $w_i$ \\
    \midrule
    $S_\text{acc}(\mathrm{TimingAcc})$ & 0.50 \\
    $S_\text{pw}(\mathrm{MAE}_\text{on})$ & 0.30 \\
    $S_\text{rhythm}$ & 0.20 \\
    \bottomrule
  \end{tabular}
\end{center}
where:
\[
  S_\text{rhythm}
    = \tfrac{1}{2}\bigl(S_\text{pw}(\mathrm{IOI\text{-}MAE})
    + S_\text{pw}(\sigma_\text{on})\bigr)
\]

\textbf{Duration category} (weight 0.15 in final score):
\begin{center}
  \begin{tabular}{@{}lr@{}}
    \toprule
    Component & Weight $w_i$ \\
    \midrule
    $S_\text{pw}(\mathrm{RelDurErr})$ & 0.60 \\
    $S_\text{pw}(|\mathrm{DurRatio}-1|)$ & 0.20 \\
    $S_\text{phrase}$ & 0.20 \\
    \bottomrule
  \end{tabular}
\end{center}
where:
\[
  S_\text{phrase}
    = \tfrac{1}{2}\bigl(S_\text{pw}(\sigma_\text{dur})
    + S_\text{pw}(\mathrm{RelDurErr})\bigr)
\]

\textbf{Lyric category} (weight 0.15 in final score):
\begin{center}
  \begin{tabular}{@{}lr@{}}
    \toprule
    Component & Weight $w_i$ \\
    \midrule
    $S_\text{acc}(\mathrm{WordAcc})$    & 0.35 \\
    $S_\text{acc}(\mathrm{OverlapAcc})$ & 0.25 \\
    $S_\text{acc}(\mathrm{LabelMatch})$ & 0.25 \\
    $S_\text{pw}(\mathrm{PhonBndErr})$  & 0.15 \\
    \bottomrule
  \end{tabular}
\end{center}

\subsection{Final Score}

The overall coaching score is a confidence-adjusted weighted average of the
four category scores.
Let $W_k$ be the nominal weight and $C_k = \min(1, N_k/3)$ be the confidence
of category $k$.

\[
  \boxed{S_\text{final}
    = \frac{\sum_k W_k\,C_k\,S_k}{\sum_k W_k\,C_k}}
  \qquad
  \text{nominal: } W_\text{pitch}=0.40,\;
    W_\text{timing}=0.30,\;
    W_\text{dur}=0.15,\;
    W_\text{lyric}=0.15
\]

Pitch receives the highest weight (0.40) because intonation accuracy is
the most perceptually salient dimension of singing quality.

% =============================================================================
\section{Interpretation (Phase 7c)}
% =============================================================================

\texttt{scoring/interpretation.py} maps the final and category scores to
human-readable feedback using a fully deterministic rule table.
No language models or probabilistic text generation are involved.

\textbf{Score thresholds:}
\begin{center}
  \begin{tabular}{@{}clp{7cm}@{}}
    \toprule
    Score range & Level & Meaning \\
    \midrule
    $\ge 90$ & \texttt{excellent}  & Performance-ready; present-ready. \\
    $\ge 75$ & \texttt{good}       & Strong performance; a few targeted corrections. \\
    $\ge 55$ & \texttt{fair}       & Foundational elements present; focused practice needed. \\
    $< 55$   & \texttt{needs\_work} & Coaching required before full-song practice. \\
    \bottomrule
  \end{tabular}
\end{center}

Strengths are emitted for categories at \texttt{excellent} or \texttt{good}
levels; weaknesses are emitted for \texttt{fair} or \texttt{needs\_work}.
All messages are pre-defined templates, e.g.:
\begin{itemize}[noitemsep]
  \item Strength (pitch, excellent): ``Strong pitch intonation and accuracy.''
  \item Weakness (timing, needs\_work): ``Significant timing inconsistencies
        detected.''
\end{itemize}

Output: \texttt{InterpretationSummary(overall\_level, strengths, weaknesses,
category\_levels)}.

% =============================================================================
\section{Configuration}
% =============================================================================

All pipeline parameters are controlled via YAML files in \texttt{configs/}.
The master file \texttt{pipeline.yaml} governs feature flags, model
checkpoints, fusion thresholds, alignment tolerances, metric tolerances, and
scoring weights.
Per-model files (\texttt{pitch.yaml}, \texttt{onset\_offset.yaml},
\texttt{phoneme.yaml}) control model-specific hyperparameters.
Table~\ref{tab:config} summarizes the most important configurable values.

\begin{table}[h]
  \centering
  \caption{Key configuration parameters and their defaults.}
  \label{tab:config}
  \renewcommand{\arraystretch}{1.2}
  \begin{tabular}{@{}lll@{}}
    \toprule
    Parameter & Default & Effect \\
    \midrule
    \texttt{onset\_threshold}   & 0.5     & Peak-pick threshold for note onset detection \\
    \texttt{offset\_threshold}  & 0.5     & Peak-pick threshold for note offset detection \\
    \texttt{min\_note\_duration\_ms} & 50 & Discard notes shorter than this \\
    \texttt{word\_gap\_ms}      & 100     & Phoneme gap that starts a new word \\
    \texttt{phrase\_gap\_s}     & 0.5     & Note gap that starts a new phrase \\
    \texttt{alignment\_tolerance\_ms} & 100 & Max onset deviation for note matching \\
    \texttt{pitch\_tolerance\_cents}  & 50  & Tolerance for pitch accuracy metric \\
    \texttt{timing\_tolerance\_ms}    & 50  & Tolerance for timing accuracy metric \\
    \texttt{phoneme\_tolerance\_ms}   & 30  & Tolerance for phoneme timing accuracy \\
    \texttt{VAD aggressiveness}       & 2   & WebRTC VAD sensitivity \\
    \texttt{median\_filter\_size}     & 5   & Frames for VAD/pitch median filter \\
    \texttt{max\_gap\_fill\_frames}   & 10  & Max frames to linearly interpolate \\
    \texttt{smoothing\_sigma}         & 1.0 & Gaussian smoothing $\sigma$ (frames) \\
    \texttt{max\_segment\_ms}         & 300 & Max phoneme segment before forced split \\
    \bottomrule
  \end{tabular}
\end{table}

% =============================================================================
\section{Data Types and Contracts}
% =============================================================================

The system is strongly typed via dataclasses defined in \texttt{utils/types.py}.
Every inter-phase boundary has a well-defined Python dataclass as its contract.
Table~\ref{tab:types} lists the principal types and the pipeline phase that
produces each.

\begin{table}[h]
  \centering
  \caption{Principal dataclass types and the phase that produces them.}
  \label{tab:types}
  \renewcommand{\arraystretch}{1.2}
  \begin{tabular}{@{}lll@{}}
    \toprule
    Type & Phase & Description \\
    \midrule
    \texttt{PhonemeSegment}                  & 3 & Predicted phoneme with time bounds \\
    \texttt{NoteEvent}                       & 5 & Predicted note with pitch stats \\
    \texttt{LyricEvent}                      & 5 & Single phoneme + word/note refs \\
    \texttt{WordEvent}                       & 5 & Group of phonemes forming a word \\
    \texttt{PhraseEvent}                     & 5 & Group of notes forming a phrase \\
    \texttt{FrameAlignedFeatures}            & 5 & All model outputs at 100\,fps \\
    \texttt{FusedPerformanceRepresentation}  & 5 & Complete predicted-side representation \\
    \texttt{ReferencePerformanceRepresentation} & 6a & Ground-truth notes, phonemes, words \\
    \texttt{AlignmentResult}                 & 6b & Matched pairs + unmatched indices \\
    \texttt{PitchMetrics}                    & 7a & Pitch metric values \\
    \texttt{TimingMetrics}                   & 7a & Timing metric values \\
    \texttt{DurationMetrics}                 & 7a & Duration metric values \\
    \texttt{LyricMetrics}                    & 7a & Lyric metric values \\
    \texttt{CategoryScore}                   & 7b & Weighted aggregate per dimension \\
    \texttt{PerformanceScoreReport}          & 7b & All category + overall scores \\
    \texttt{InterpretationSummary}           & 7c & Level labels + feedback strings \\
    \bottomrule
  \end{tabular}
\end{table}

All types implement a \texttt{to\_dict()} method for JSON serialization,
enabling structured logging and offline analysis.

% =============================================================================
\section{Pilot Results}
% =============================================================================

A pilot evaluation was conducted on a 7-note excerpt from the song
``North Wind Meets the Sea'' using the torchcrepe backend and a
tolerance of $\tau = 50$\,¢.

\[
  \mathrm{NotePitchAcc}_{50} = 6/7 \approx 85.7\%
\]

\begin{table}[h]
  \centering
  \caption{Per-note pitch results on the pilot excerpt.}
  \renewcommand{\arraystretch}{1.2}
  \begin{tabular}{@{}clccccc@{}}
    \toprule
    Note & Lyric & Ref (MIDI) & Ref (Hz) & Det.\ (MIDI) & $\bar{c}_n$ (¢) & Correct \\
    \midrule
    1 & \textit{where} & C4 (60.0)  & 261.6 & 59.74 & $-26.1$  & $\checkmark$ \\
    2 & \textit{the}   & D4 (62.0)  & 293.7 & 60.47 & $-152.8$ & $\times$     \\
    3 & \textit{north} & E4 (64.0)  & 329.6 & 64.08 & $+8.2$   & $\checkmark$ \\
    4 & \textit{wind}  & B$\flat$3 (58.0) & 233.1 & 57.63 & $-36.8$ & $\checkmark$ \\
    5 & \textit{meets} & C4 (60.0)  & 261.6 & 60.32 & $+32.5$  & $\checkmark$ \\
    6 & \textit{the}   & A3 (57.0)  & 220.0 & 57.44 & $+44.4$  & $\checkmark$ \\
    7 & \textit{sea}   & D4 (62.0)  & 293.7 & 62.24 & $+23.7$  & $\checkmark$ \\
    \bottomrule
  \end{tabular}
\end{table}

Note 2 is the only failure at $-152.8$\,¢ (approximately 1.5 semitones flat).
Its duration is only 178\,ms with 18 voiced frames, which is insufficient
data for a reliable median F0 estimate.
All six correctly identified notes fall within $\pm 45$\,¢.
There is a mild systematic bias: the model tends to be flat on longer
sustained notes (notes 1, 4) and slightly sharp on shorter notes (notes 3,
5, 6, 7), which warrants further investigation as the evaluation dataset
expands.

% =============================================================================
\section{Architectural Design Principles}
% =============================================================================

The following design choices distinguish \texttt{VocalCoach\_Kim} from
end-to-end score-predicting systems.

\textbf{Neural networks as feature extractors, not oracles.}
Each model predicts a physical, interpretable quantity (F0, note boundaries,
phoneme boundaries) rather than a quality score.
This means errors are diagnosable: a wrong pitch score implies a wrong F0
estimate, which can be inspected in the pitch JSON log.

\textbf{Deterministic scoring.}
Given the same audio and reference files, the system always produces the
same score.
There is no stochastic decoding, no sampling, and no LLM call anywhere in
the scoring path.

\textbf{Confidence-aware aggregation.}
Scores are down-weighted when very few notes or phonemes are matched.
A 3-note evaluation does not carry the same statistical weight as a
100-note evaluation; the confidence factor $c_i = \min(1, N_i/3)$ reflects
this.

\textbf{Separation of prediction and evaluation.}
Phases 1--5 produce the predicted-side representation independently of any
reference data.
Phases 6--7 introduce reference data.
This means the prediction pipeline can be run in a ``reference-free'' mode
(e.g., for live monitoring) without modification.

\textbf{Lazy model loading.}
Models are loaded on demand via a central \texttt{ModelRegistry}, not at
import time.
This keeps startup latency low when only a subset of the pipeline is needed.

\textbf{Configurable normalization.}
The breakpoints for all piecewise-linear scoring curves are stored in YAML
and can be tuned per deployment (e.g., stricter thresholds for a conservatory
context, looser thresholds for a beginner coaching app) without touching
source code.

% =============================================================================
\section*{Acknowledgements}
% =============================================================================

This system was developed at the IMAPLE Lab, Drexel University.
The pitch estimation backbone uses the torchcrepe implementation of the CREPE
model.
Onset/offset and phoneme detection use pre-trained models from the
\texttt{facebook/wav2vec2} family.
Reference data parsing uses the \texttt{music21} and \texttt{praatio}
libraries.

\end{document}
